% Source printed pp. 508-510 / PDF pp. 516-518.
\subsection{Abel 簇}\label{src:10.1.5-subsection}
\newterm{主极化 Abel 簇}{principally polarized abelian variety}的模空间 $A_g$ 将椭圆曲线的模空间推广到高维情形. 在过去两个世纪中, $A_g$ 在代数几何和算术几何中发挥了重要作用. 它还通过 Torelli 映射 $M_g\to A_g$ 与 $M_g$ 深刻相连. 在 Riemann 思想的基础上, Baily, Igusa, Siegel, Satake 等人在二十世纪中叶发展了 $A_g$ 的解析理论; 随后 Chai, Faltings, Mumford, Oort 及许多其他数学家建立了代数理论.

\textbf{紧复环面的模空间.} 一个 $g$ 维\newterm{复环面}{complex torus} $X$ 是带有群结构的 $g$ 维连通紧复流形. 将 $\mathbb C^g$ 与 $T_{X,0}$ 等同后, 指数映射 $\exp:\mathbb C^g\to X$ 的核 $\Lambda$ 是格 $\Lambda=\mathbb Z^{2g}\subseteq\mathbb C^g$, 且 $X$ 可与商 $\mathbb C^g/\Lambda$ 等同. 定义满秩格 $\Lambda=\im(\lambda)$ 的群同态 $\lambda:\mathbb Z^{2g}\to\mathbb C^g$ 所成的空间同构于 $\mathrm{GL}_{2g}(\mathbb R)$. 两个复环面 $\mathbb C^g/\im(\lambda)$ 和 $\mathbb C^g/\im(\lambda')$ 同构, 当且仅当存在矩阵 $\sigma\in\mathrm{GL}_g(\mathbb C)$ 和 $\sigma'\in\mathrm{GL}_{2g}(\mathbb Z)$, 满足 $\lambda'=\sigma\circ\lambda\circ\sigma'$. 换言之, 复环面的模空间是双商
\[
\mathcal M_{\mathrm{cpx\ tori}}=\mathrm{GL}_g(\mathbb C)\backslash\mathrm{GL}_{2g}(\mathbb R)/\mathrm{GL}_{2g}(\mathbb Z),
\]
其复维数为 $g^2$. 然而, $\mathrm{GL}_{2g}(\mathbb Z)$ 在 $\mathrm{GL}_g(\mathbb C)\backslash\mathrm{GL}_{2g}(\mathbb R)$ 上的作用并非正常不连续, 商也不是 Hausdorff 的.

\textbf{\texorpdfstring{$\mathbb C$}{C} 上 Abel 簇的模空间.} 复环面 $\mathbb C^g/\Lambda$ 是代数的, 当且仅当存在正定 Hermite 形式 $H:\mathbb C^g\times\mathbb C^g\to\mathbb C$, 使其虚部 $\alpha:=\operatorname{im}(H)$ 在 $\Lambda$ 上取整数值 \cite{Rie57}; 此时称 $\alpha$ 为 \newterm{Riemann 形式}{Riemann form}. $\mathbb C^g/\Lambda$ 的代数性也可由亚纯函数域的超越次数为 $g$ 来刻画; 见 \cite[第 35 页推论]{Mum70a}. 并非每个复环面都是代数的; 例如, Siegel 给出了一个没有非常值亚纯函数的 $2$ 维复环面 \cite[第 104 页]{Sie50}. 事实上, 对几乎所有格 $\Lambda$, 复环面 $\mathbb C^g/\Lambda$ 都不是代数的; 见 \cite[第 37 页引理]{Mum70a}.

固定 $\mathbb Z^{2g}$ 的辛基 $(\gamma_i)$ 及其上的交替形式 $\alpha$, 并要求 $\alpha$ 幺模, 即 $\det(\alpha(\gamma_i,\gamma_j))=1$. 由 Riemann 形式 $\alpha$ 极化的代数复环面 $\mathbb C^g/\Lambda$ 由辛矩阵 $\lambda\in\mathrm{Sp}_{2g}(\mathbb R)$ 确定, 该矩阵将 $\Lambda$ 与 $\lambda(\mathbb Z^{2g})$ 等同. 两个矩阵 $\lambda,\lambda'\in\mathrm{Sp}_{2g}(\mathbb R)$ 定义同构的 Abel 簇, 当且仅当存在 $\sigma\in U_g(\mathbb C)$ 和 $\gamma\in\mathrm{Sp}_{2g}(\mathbb Z)$, 满足 $\lambda'=\sigma\circ\lambda\circ\gamma$. 因此 $g$ 维 Abel 簇的模空间可表为双商
\[
A_g=U_g(\mathbb C)\backslash\mathrm{Sp}_{2g}(\mathbb R)/\mathrm{Sp}_{2g}(\mathbb Z),
\]
其复维数为 $g(g+1)/2$. 商 $U_g(\mathbb C)\backslash\mathrm{Sp}_{2g}(\mathbb R)$ 与 \newterm{Siegel 上半空间}{Siegel upper-half plane}
\[
\mathcal H_g=\{\tau\in\mathrm{Mat}_{g\times g}(\mathbb C):\tau=\tau^t,\ \operatorname{Im}(\tau)>0\}
\]
等同, 其中元素 $\tau\in\mathcal H_g$ 定义代数复环面 $\mathbb C^g/(\mathbb Z^g+\tau\mathbb Z^g)$. $\mathrm{Sp}_{2g}(\mathbb Z)$ 在 $\mathcal H_d$ 上的作用正常不连续, 因此模空间是复解析空间. 见 \cite[\S II.4]{Mum07} 和 \cite[\S10]{Gee08}.

\textbf{\texorpdfstring{$\mathbb Z$}{Z} 上 Abel 簇的模空间.} 抽象 Abel 簇的模叠不是代数叠, 因为存在非有效的形式形变; 见注记~\bookref{C.5.9}. 为了解决这个问题, 我们研究带有主极化的 Abel 簇, 主极化将在下面定义.

概形 $S$ 上的 \newterm{Abel 概形}{abelian scheme}是具有连通几何纤维的光滑固有群概形 $\pi:A\to S$. 群律 $m:A\times_S A\to A$ 必然交换. Picard 函子 $\underline{\mathrm{Pic}}_{A/S}$ 可表 (命题~\bookref{7.4.62}), 开子群概形 $\widehat A:=\underline{\mathrm{Pic}}_{A/S}^\tau$ 参数化某个幂代数等价于零的线丛, 它也是 $S$ 上的 Abel 概形, 称为\newterm{对偶 Abel 概形}{dual abelian scheme}. $A$ 上的线丛 $L$ 诱导 $A\times_S A$ 上的线丛 $m^*(L)\otimes p_1^*(L)^\vee\otimes p_2^*(L)^\vee$, 进而诱导映射 $\psi_L:A\to\widehat A$. 当 $L$ 丰沛时, $\psi_L$ 是同源, 称 $\psi_L$ 为\newterm{极化}{polarization}. 若 $\psi_L$ 还是同构, 则称为\newterm{主极化}{principal polarization}. 在代数闭域上, 丰沛线丛 $L$ 定义主极化当且仅当 $h^0(A,L)=1$. 在 $\mathbb C$ 上, 主极化对应于幺模 Riemann 形式. 与曲线的情形 (命题~\bookref{6.1.16}) 类似, 若 $(A\to S,\psi_L)$ 是主极化 Abel 概形, 则当 $k\geq3$ 时 $L^{\otimes k}$ 相对甚丰沛 \cite[定理 6.13]{GIT}; 这可用于建立代数性和有界性.

主极化 Abel 簇的叠 $\mathcal A_g$ 的对象是带有主极化 $\psi_L$ 的 Abel 概形 $A\to S$. 它是 $\mathbb Z$ 上维数为 $g(g+1)/2$ 的分离有限型 Deligne--Mumford 叠, 具有拟射影粗模空间 $A_g$ \cite[定理 7.10]{GIT}.\footnote{与对 $M_g$ 的论证一样, Mumford 的论证只证明了对每个素数 $p$, 在 $\Spec\mathbb Z\setminus\{p\}$ 上的拟射影性 (见注记~\bookref{5.4.7}).}

模空间 $A_g$ 不紧, 并有若干不同的紧化 --- Bailey--Borel 紧化, Satake 紧化, 以及 Igusa 和 Mumford 的环面紧化 --- 各有优缺点. 环面紧化包括第一和第二 Voronoi 紧化. Faltings 证明了这些紧化在 $\mathbb Z$ 上存在 \cite{FC90}. 这些紧化以函子方式定义, 而 Alexeev 证明了第二 Voronoi 紧化作为某个模性紧化的一个不可约分支出现 \cite{Ale02}. Olsson 随后利用对数几何给出了这个不可约分支的函子描述 \cite{Ols08}. 关于 Abel 簇的模空间及其紧化的综述, 推荐 \cite{GO99,Gru09}.
